% angle_repose: a block on the verge of sliding at the angle of repose (Ch 9, fig:repose)
% Drawn with theta_r = 30 degrees (so mu_s = tan 30 = 0.577) and the forces to scale:
%   mg = 2.4 cm,  N = mg cos(theta_r),  limiting friction mu_s N = mg sin(theta_r).
% N and mu_s N act at the contact; their resultant R (the total contact force) is
% vertical and equal to mg, so it leans from the normal by lambda = theta_r: the
% angle of friction and the angle of repose are the same angle.
\begin{tikzpicture}[x=1cm,y=1cm]
  \def\th{30}
  \def\Lb{8.6}                                  % base of the incline
  \def\W{2.4}                                   % mg, in cm
  \def\s{4.8}                                   % block's contact point, along the slope
  \pgfmathsetmacro\Hh{\Lb*tan(\th)}
  % incline and the floor
  \path[ground,draw=none] (-0.35,-0.14) rectangle (\Lb+0.35,0);
  \fill[cSlope] (0,0) -- (\Lb,0) -- (\Lb,\Hh) -- cycle;
  \draw[outline] (-0.35,0) -- (\Lb+0.35,0);
  \draw[outline] (0,0) -- (\Lb,\Hh) -- (\Lb,0);
  % theta_r at the foot of the incline
  \draw[angle arc] (1.15,0) arc[start angle=0,end angle=\th,radius=1.15];
  \node at (\th/2:1.5) {$\theta_r$};
  % the block, in the slope's own frame
  \coordinate (P) at (\th:\s);                  % contact point (centre of the base)
  \begin{scope}[shift={(P)},rotate=\th]
    \draw[body] (-0.75,0) rectangle (0.75,0.95);
    \coordinate (C) at (0,0.475);               % centre of mass
  \end{scope}
  % contact forces at P, to scale, and their resultant R (dashed: a resultant, not a
  % separate force)
  \pgfmathsetmacro\Nl{\W*cos(\th)}
  \pgfmathsetmacro\fl{\W*sin(\th)}
  \coordinate (Nt) at ($(P)+(90+\th:\Nl)$);
  \coordinate (ft) at ($(P)+(\th:\fl)$);
  \coordinate (Rt) at ($(P)+(0,\W)$);
  \draw[parallelogram] (Nt) -- (Rt) -- (ft);
  \draw[force,resultant] (P) -- (Rt) node[anchor=south] {$R$};
  \draw[force] (P) -- (Nt) node[anchor=south east,inner sep=1pt] {$N$};
  \draw[applied] (P) -- (ft) node[anchor=north west,inner sep=2pt] {$\mu_s N$};
  % lambda: how far the contact force leans from the normal
  \draw[angle arc] ($(P)+(90:1.5)$) arc[start angle=90,end angle=90+\th,radius=1.5];
  \node at ($(P)+(90+\th/2:1.8)$) {$\lambda$};
  % weight at the centre of mass
  \draw[force] (C) -- ++(0,-\W) node[anchor=west,inner sep=2pt,pos=0.9] {$mg$};
  \node[contact point] at (P) {};
\end{tikzpicture}
